% Living methods notes for the ladder MPS+MF project.
% Keep this file synchronized with the modular implementation under ../src/.
% Scientific run logs and campaign decisions belong in RUN_LOG.md, not here.
% Compile from this directory with: pdflatex METHODS_NOTES.tex (twice).

\documentclass[11pt]{article}

\usepackage[T1]{fontenc}
\usepackage{amsmath,amssymb,bm}
\usepackage{booktabs}
\usepackage{geometry}
\usepackage{hyperref}
\usepackage{tabularx}
\usepackage{xcolor}

\geometry{margin=1in}
\hypersetup{
  colorlinks=true,
  linkcolor=blue!55!black,
  urlcolor=blue!55!black,
  citecolor=blue!55!black
}
\setlength{\parindent}{0pt}
\setlength{\parskip}{0.55em}
\renewcommand{\arraystretch}{1.15}
\setcounter{tocdepth}{2}

\newcommand{\Hlad}{H_{\mathrm{lad}}}
\newcommand{\Heff}{H_{\mathrm{eff}}}
\newcommand{\Ep}{E_{\mathrm p}}
\newcommand{\tp}{t_{\perp}}
\newcommand{\Evar}{E_{\mathrm{var}}}
\newcommand{\mean}[1]{\left\langle #1\right\rangle}
\newcommand{\dd}{\mathrm d}
\newcommand{\hc}{\mathrm{H.c.}}

\title{Ladder MPS+Mean-Field Project: Living Methods Notes}
\author{Kevin Wang and contributors}
\date{Initialized 1 September 2026}

\begin{document}

\maketitle

\begin{abstract}
These notes summarize the methods currently implemented in the modular
\texttt{ladder\_mps\_mft} workflow.  The calculation combines finite-system
matrix-product-state (MPS) ground states for an interacting two-leg ladder with
a perturbative transverse mean-field (MF) map, explicit density targeting,
fixed-point and physical-orbit validation, and a zero-temperature variational
energy comparison.  This is an implementation-oriented living document, not a
claim of thermodynamic convergence.  Unless stated otherwise, ``accepted''
means numerically self-consistent at the simulated length and bond dimension.
\end{abstract}

\tableofcontents

\section{Scope and notation}

The authoritative implementation is the Julia package in
\path{ladder_mps_mft/src}.  The monolithic scripts at the repository root are
retained for legacy provenance but are not the definition of the current
method.

The simulated object is one open two-leg ladder with \(L\) rungs and \(2L\)
physical sites.  A site is denoted \(a=(i,\ell)\), with rung
\(i=1,\ldots,L\), leg \(\ell\in\{0,1\}\), and spin
\(\sigma\in\{\uparrow,\downarrow\}\).  The MPS ordering is rung-major,
\((1,0),(1,1),(2,0),(2,1),\ldots\).  The stored density is per physical site,
\begin{equation}
  n = \frac{N}{2L}, \qquad N_{\star}=2L n_{\star}.
\end{equation}
All present MF fields and measured correlators are restricted to their real
parts.

\section{Isolated-ladder Hamiltonian}

The bare ladder Hamiltonian is the extended Hubbard model
\begin{align}
\Hlad={}&-t\sum_{i=1}^{L-1}\sum_{\ell,\sigma}
 \left(c^{\dagger}_{i\ell\sigma}c_{i+1,\ell\sigma}+\hc\right)
 -t_0\sum_{i=1}^{L}\sum_{\sigma}
 \left(c^{\dagger}_{i0\sigma}c_{i1\sigma}+\hc\right) \nonumber\\
&+U\sum_{i,\ell}n_{i\ell\uparrow}n_{i\ell\downarrow}
 +V\sum_{\langle a,b\rangle_{\mathrm{lad}}}n_a n_b .
\label{eq:bare-hamiltonian}
\end{align}
Here \(t\) is the leg hopping, \(t_0\) the rung hopping, and the \(V\) sum
contains nearest-neighbor leg and rung bonds.  Open boundary conditions are
used along the ladder.  The inter-ladder hopping \(\tp\) is not inserted as an
explicit third spatial direction; it enters the effective MF map at second
order through \(\tp^2/|\Ep|\).

\subsection{Pair-binding scale and registry}

The hole pair-binding convention used by the registry is
\begin{equation}
 \Ep^{(h)}=E(N-2,2S_z=0)+E(N,2S_z=0)
             -2E(N-1,2S_z=1).
\label{eq:pair-binding}
\end{equation}
A negative value denotes a bound pair.  The MF denominator is positive,
\begin{equation}
  g=\frac{\tp^2}{|\Ep|}, \qquad p=2g.
\label{eq:g-definition}
\end{equation}
By default the code refuses an unbound or zero \(\Ep\).  It first seeks an
exact registry match in \((L,U,V,t_0,n)\), choosing the largest available bond
dimension and then the smallest recorded relative discrepancy.  Optional
interpolation is linear in the \emph{signed} \(\Ep\), only between bracketing
\(t_0\) values with identical \((L,U,V,n)\).  Extrapolation and interpolation
across a sign change are forbidden.  The selected row, endpoints, weights,
and registry hash are stored with every result.

The inexpensive check \(\tp<|\Ep|\) is necessary but not a complete
weak-coupling test.  Where available, \(\tp\) is also compared with the
isolated-ladder spin and charge gaps defined in Sec.~\ref{sec:diagnostics}.

\section{Transverse mean-field construction}
\label{sec:mean-field}

\subsection{Fields, correlators, and effective Hamiltonian}

At SCF step \(k\), collect the applied fields into
\(x_k=(\alpha,\beta,\nu)_k\):
\begin{itemize}
  \item \(\alpha_{ii'}^{\ell\ell'}\): anomalous singlet-pair field;
  \item \(\beta_{\sigma,ii'}^{\ell\ell'}\): spin-resolved normal exchange
        field, with the on-site entry removed; and
  \item \(\nu_{\sigma,i\ell}\): spin-resolved Hartree field, stored in code as
        \texttt{mu\_cdw} and distinct from the global chemical potential.
\end{itemize}
Only pairs with \(|i-i'|\le r_{\mathrm{range}}\) are retained.  Define the real
correlators
\begin{align}
 F_{i'\ell',i\ell}&=\operatorname{Re}
   \mean{c_{i'\ell'\uparrow}c_{i\ell\downarrow}},\\
 X^{\sigma}_{a,b}&=\operatorname{Re}
   \mean{c^{\dagger}_{a\sigma}c_{b\sigma}},
 \qquad n^{\sigma}_{a}=\mean{n_{a\sigma}}.
\end{align}
The effective single-ladder problem is
\begin{equation}
  \Heff[x_k,\mu_k]=\Hlad-\mu_k N+H_{\mathrm{lin}}[x_k],
\label{eq:effective-hamiltonian}
\end{equation}
with
\begin{align}
H_{\mathrm{lin}}={}&
-\sum_{ii'\ell\ell'}\alpha_{ii'}^{\ell\ell'}
 \left(c_{i'\ell'\uparrow}c_{i\ell\downarrow}
       +c^{\dagger}_{i\ell\downarrow}c^{\dagger}_{i'\ell'\uparrow}\right)
\nonumber\\
&+\sum_{ii'\ell\ell'\sigma}\beta_{\sigma,ii'}^{\ell\ell'}
 c^{\dagger}_{i\ell\sigma}c_{i'\ell'\sigma}
 +\sum_{i\ell\sigma}\nu_{\sigma,i\ell}n_{i\ell\sigma}.
\label{eq:linear-fields}
\end{align}

\subsection{Geometry-dependent map}

Solving Eq.~\eqref{eq:effective-hamiltonian} and measuring \(F,X,n\) defines
the raw map
\begin{equation}
  y_k=f(x_k), \qquad x_{k+1}=y_k
\label{eq:raw-map}
\end{equation}
during an unmixed probe.  The geometry labels are
\texttt{cubic\_frustrated}, \texttt{cubic\_unfrustrated}, and
\texttt{square}.  For the same-leg anomalous fields, let
\(\bm F=(F_{i'0,i0},F_{i'1,i1})^{\mathsf T}\).  Their exact implemented map is
\begin{equation}
 \begin{pmatrix}\alpha^{00}_{ii'}\\\alpha^{11}_{ii'}\end{pmatrix}
 =p M_G\bm F,
\end{equation}
where
\begin{center}
\begin{tabular}{lll}
\toprule
Geometry \(G\) & \(M_G\) & Cross-leg anomalous output
\\ \midrule
frustrated cubic
& \(\left(\begin{smallmatrix}2&1\\1&2\end{smallmatrix}\right)\)
& \((\alpha^{10},\alpha^{01})=2p(F_{i'0,i1},F_{i'1,i0})\)
\\
unfrustrated cubic
& \(\left(\begin{smallmatrix}0&3\\3&0\end{smallmatrix}\right)\)
& \((\alpha^{10},\alpha^{01})=2p(F_{i'1,i0},F_{i'0,i1})\)
\\
square
& \(\left(\begin{smallmatrix}0&1\\1&0\end{smallmatrix}\right)\)
& \((\alpha^{10},\alpha^{01})=(0,0)\)
\\ \bottomrule
\end{tabular}
\end{center}
The normal fields \(\beta_\sigma\) use the same integer coefficient pattern
with \(F\) replaced by the corresponding \(X^\sigma\) entries and the normal
operator orientation in Eq.~\eqref{eq:linear-fields}; their diagonal on-site
entries are set to zero.  This component-level map is implemented in
\path{src/MeanField.jl}; changes to it require a simultaneous audit of the
variational constants in Sec.~\ref{sec:variational}.

For each rung and spin, the Hartree update is
\begin{equation}
 \begin{pmatrix}\nu_{\sigma,i0}\\\nu_{\sigma,i1}\end{pmatrix}
 =K_G\left[
 \begin{pmatrix}n^\sigma_{i0}\\n^\sigma_{i1}\end{pmatrix}
 -\frac12\begin{pmatrix}1\\1\end{pmatrix}\right],
\label{eq:density-kernel}
\end{equation}
with
\begin{center}
\begin{tabular}{lcc}
\toprule
Geometry \(G\) & \(K_G\) & leg-even / leg-odd eigenvalues \\
\midrule
frustrated cubic
& \(2g\left(\begin{smallmatrix}2&1\\1&2\end{smallmatrix}\right)\)
& \(6g,\;2g\) \\
unfrustrated cubic
& \(6g\left(\begin{smallmatrix}0&1\\1&0\end{smallmatrix}\right)\)
& \(6g,\;-6g\) \\
square
& \(2g\left(\begin{smallmatrix}0&1\\1&0\end{smallmatrix}\right)\)
& \(2g,\;-2g\) \\
\bottomrule
\end{tabular}
\end{center}
Negative map eigenvalues naturally motivate a two-sublattice transverse
interpretation and can appear in the one-ladder iteration as a period-two
orbit.  Numerical recurrence alone, however, does not identify its physical
order parameter.

\section{MPS ground-state solve and density targeting}

\subsection{Finite-system DMRG}

For fixed \(x_k\) and \(\mu\), the ground state of \(\Heff\) is approximated by
finite-system DMRG using ITensorMPS.  A cold start uses a randomized product
state with particle count nearest \(2Ln_\star\); subsequent MF and chemical
potential solves warm-start from the previous MPS.  The observer records each
sweep energy, maximum discarded weight, and maximum link dimension.  It may
terminate only after reaching the final configured bond dimension with zero
noise and satisfying the sweep-energy tolerance, or when the wall-time
deadline is reached.

Production GPU calculations use dense tensors in Float64, with \(S_z\) and
fermion-parity quantum-number block structure disabled for performance.  This
changes the tensor representation, not the Hamiltonian.  Separate fixed-sector
gap calculations use QN-conserving CPU MPS states.

\subsection{Safeguarded chemical-potential search}

Each MF update contains an inner solve of
\begin{equation}
  h(\mu)=n(\mu;x_k)-n_\star=0.
\end{equation}
The algorithm first evaluates the previous \(\mu\).  If a positive carried
slope \(\kappa\simeq\dd n/\dd\mu\) is available, it tries a bounded Newton
predictor.  Otherwise, or if that predictor does not bracket the target, it
expands a bracket geometrically in the density-correct direction.  Within a
bracket it uses a safeguarded secant step, rejecting points within ten percent
of either endpoint in favor of bisection.  Every trial reuses the nearest MPS
and a small warm-start noise schedule.  The best-density state is returned on
an iteration limit, but it cannot pass final acceptance unless the configured
density gate is met.

\section{Self-consistency, mixing, and recurrence}
\label{sec:scf}

One outer iteration is:
\begin{enumerate}
  \item solve \(\Heff[x_k,\mu_k]\) by density-targeted DMRG;
  \item measure \(y_k=f(x_k)\) and the associated correlators;
  \item evaluate field residuals and the variational-energy identities;
  \item classify the recent history as a fixed point, physical orbit,
        candidate, or incomplete numerical trajectory; and
  \item if nonterminal, choose \(x_{k+1}\) by the raw map or an accelerator.
\end{enumerate}

The field residuals use the complete stored field vector,
\begin{equation}
 \epsilon_{\infty}=\max_j|(y_k-x_k)_j|,\qquad
 \epsilon_2=\frac{\|y_k-x_k\|_2}
 {\max(\|y_k\|_2,\|x_k\|_2,\epsilon_{\mathrm{mach}})}.
\label{eq:field-residuals}
\end{equation}

\subsection{Unmixed probe and physical orbits}

The production-shaped workflow begins with complete raw updates
\(x_{k+1}=f(x_k)\), before damping or Anderson acceleration.  A period-\(q\)
recurrence requires every phase to close over the configured number of repeat
links.  With three repeat links, a period-two test needs eight contiguous raw
records.  Acceptance additionally requires:
\begin{itemize}
  \item raw-map closure \(x_k\simeq y_{k-1}\) for every stored link;
  \item target density and per-phase energy recurrence;
  \item the Hamiltonian and effective-eigenvalue consistency checks in
        Sec.~\ref{sec:variational}; and
  \item explicit inclusion of \(q\) in the configured accepted-period list.
\end{itemize}
For \(q=2\), successive step vectors must reverse direction: the worst recent
cosine must be below its negative threshold, and
\(\|y_k-y_{k-2}\|/\|y_k-y_{k-1}\|\) must remain below its configured bound.
This rejects a slowly drifting trajectory that happens to be close at lag two.

A mixer-dependent recurrence is only a candidate.  It is archived, mixer
history is cleared, and a fresh raw-map probe is required before it may be
accepted.  Orbit phases are stored separately and never replaced by their
field average.  The physical period-two construction follows the logic of
Ref.~\cite{bollmark2025}; naming an accepted orbit ``CDW'' still requires a
nonzero bulk density contrast, the expected spatial pattern, and stability
with \(L\), bond dimension, and bulk window.  The stored contrast is the mean
absolute density difference between neighboring orbit phases over a centered
bulk fraction of the ladder (one half by default).

\subsection{Fixed-point gate and slow modes}

A period-one solution must satisfy both the raw residual gate in
Eq.~\eqref{eq:field-residuals} and a slow-mode extrapolation.  For residual
vectors \(r_k=y_k-x_k\),
\begin{equation}
 \lambda_k=\frac{r_k\cdot r_{k-1}}{\|r_{k-1}\|_2^2}.
\end{equation}
When consecutive residuals are sufficiently parallel, the gated residual is
inflated by
\begin{equation}
  s_k=\max\left(1,\frac{1}{1-\lambda_k}\right).
\end{equation}
If \(\lambda_k\ge1\), the factor is infinite.  This prevents a nearly marginal
mode from being accepted merely because its instantaneous residual crosses a
local threshold.  The recent stability window must also pass the density gate;
the latest target-density-corrected energy change and the current energy
identities must pass their respective gates.

\subsection{Linear and Anderson acceleration}

If the initial raw probe finds no accepted solution, the code uses adaptive
linear or Anderson mixing.  Linear mixing is
\begin{equation}
 x_{k+1}=x_k+\eta_k\bigl[f(x_k)-x_k\bigr].
\end{equation}
For Anderson mixing, recent residuals \(r_j=f_j-x_j\) are combined with
coefficients \(c_j\), constrained by \(\sum_jc_j=1\), by solving the regularized
Gram/KKT system.  The update is
\begin{equation}
 x_{k+1}=(1-\eta_k)\sum_jc_jx_j+\eta_k\sum_jc_jf_j.
\end{equation}
The damping is halved after substantial residual growth and increased
conservatively after improvement, within configured minimum and maximum
bounds.  Mixing accelerates a fixed-point search; it does not certify a
physical orbit.

\subsection{Initial-condition protocol}

Independent branches deliberately sample multiple basins.  The currently
implemented seed families are summarized below.
\begin{center}
\begin{tabularx}{\textwidth}{p{0.34\textwidth}X}
\toprule
Seed family & Purpose and construction \\
\midrule
Legacy pairing/SDW/CDW
& Compatibility protocol: random retained real \(\alpha\), or deterministic
  staggered Hartree fields.  The channels do not have equal seed norms. \\
Matched mode
& Uses a declared rung profile
  \(g_m(i)=\cos\{\pi[m(i-1)/(L-1)+\phi_\pi]\}\), mean-controlled for nonzero \(m\),
  and normalizes the complete field vector so
  \(\|x\|_2/\sqrt{2L}=A_{\mathrm{seed}}\). \\
Pairing form factors
& On-site \(s\), rung \(s\), leg \(s\), extended \(s\), and \(d\)-wave
  templates, optionally with a declared longitudinal mode. \\
Stripe
& Combines a longitudinally and transversely staggered SDW envelope \(m\)
  with a leg-even charge second harmonic \(2m\). \\
Stripe+pairing
& Adds a separately normalized pairing component to test unrestricted
  coexistence; a pure stripe can remain trapped in the invariant
  \(\alpha=0\) subspace. \\
Legacy-like pairing
& Translation-invariant random relative-bond \(\alpha\) coefficients, with
  \(\beta=\nu=0\), reproducing the geometry of the old fresh-run source at the
  matched total norm. \\
\bottomrule
\end{tabularx}
\end{center}
Different seeds are initial conditions, not phase labels.  All accepted
distinct basins are retained for the variational comparison.

\section{Zero-temperature variational functional}
\label{sec:variational}

The DMRG eigenvalue \(E_{\mathrm{eig}}\) of \(\Heff\) is not directly
comparable between MF branches: it includes \(-\mu N\) and only the linearized
transverse terms.  Let \(L_{\mathrm{lin}}=\mean{H_{\mathrm{lin}}}\).  Channel
by channel,
\begin{align}
 L_{\mathrm{pair}}&=-2\sum_a\alpha_a F_a,
 &E_{\perp,\mathrm{pair}}&=\tfrac12L_{\mathrm{pair}},\\
 L_{\mathrm{ex}}&=\sum_a\beta_a X_a,
 &E_{\perp,\mathrm{ex}}&=\tfrac12L_{\mathrm{ex}},\\
 L_{\mathrm{dens}}&=\sum_a\nu_a n_a,
 &E_{\perp,\mathrm{dens}}&=\tfrac12\sum_a\nu_a(n_a-\tfrac12).
\end{align}
The applied partner fields are used in these contractions.  At a fixed point
they equal the outgoing fields; on an orbit they belong to the preceding
phase and must remain distinct.

With \(E_\perp\) the sum of the three quadratic MF contributions, the
double-counting constant is
\begin{equation}
 C_{\mathrm{dc}}=E_\perp-L_{\mathrm{lin}}.
\end{equation}
The reconstructed canonical functional is
 \begin{equation}
 E_{\mathrm{var}}^{\mathrm{rec}}
 =\mean{\Heff}+\mu\mean N+C_{\mathrm{dc}},
\end{equation}
and the independently evaluated direct form is
\begin{equation}
 E_{\mathrm{var}}^{\mathrm{dir}}=\mean{\Hlad}+E_\perp.
\label{eq:direct-functional}
\end{equation}
The code stores Eq.~\eqref{eq:direct-functional} as the canonical variational
energy and requires agreement with the reconstructed form.  It also checks
\begin{equation}
 \mean{\Hlad}=\mean{\Heff}+\mu\mean N-L_{\mathrm{lin}}
\label{eq:hamiltonian-identity}
\end{equation}
and separately compares \(E_{\mathrm{eig}}\) with \(\mean{\Heff}\).

Small residual density errors are placed on the common target-\(N\) tangent
plane using
\begin{equation}
 E_{N_\star}=E_{\mathrm{var}}^{\mathrm{dir}}
              +\mu\bigl(N_\star-\mean N\bigr).
\label{eq:density-correction}
\end{equation}
This leading correction is used for energy-stability gates and branch ranking;
it is not a substitute for tighter density convergence.  For a physical
orbit, canonical and corrected energies are computed phase by phase with the
correct partner field, then averaged arithmetically.

Only accepted fixed points or validated unmixed orbits may be ranked.  The
model, geometry, numerical settings, implementation, and \(\Ep\)-registry
fingerprints must match.  Cross-geometry energies are not ranked because the
geometries define different effective Hamiltonians, and a possible
field-independent perturbative offset has not been restored.

\section{Observables and validation diagnostics}
\label{sec:diagnostics}

\subsection{Structure factors and \(K_\rho\)}

For \(O_i=n_i\) or \(S_i^z\), the connected structure factor is
\begin{equation}
 S_O(q_x,q_y)=\frac{1}{2L}\sum_{a,b}
 e^{\mathrm{i}[q_x(i_a-i_b)+q_y(\ell_a-\ell_b)]}
 \left(\mean{O_aO_b}-\mean{O_a}\mean{O_b}\right),
\end{equation}
with \(q_x=2\pi m/L\) and \(q_y\in\{0,\pi\}\).  The dominant charge and spin
wavevectors are stored.  The reference \(q_x=\pi n\) is a tracking diagnostic,
not a universal two-leg-ladder phase criterion.

The first up to three nonzero \(q_y=0\) charge modes are fit through the origin
to obtain \(s=\dd S_c/\dd q\).  Both normalizations are reported,
\begin{equation}
 K_{\rho,\mathrm{site}}=\pi s,
 \qquad K_{\rho,\mathrm{rung}}=2\pi s.
\end{equation}

\subsection{Entanglement and pair correlations}

For every MPS bond, Schmidt probabilities \(p_a\) give
\begin{equation}
 S=-\sum_ap_a\log p_a,
 \qquad S_2=-\log\sum_ap_a^2.
\end{equation}
The central-charge diagnostic uses inter-rung cuts away from the open edges and
fits
\begin{equation}
 S(\ell)=\frac{c}{6}\log\left[\frac{2L}{\pi}
              \sin\left(\frac{\pi\ell}{L}\right)\right]+s_0.
\end{equation}
The fit window, number of points, and \(R^2\) are stored.  Finite-size,
finite-entanglement, and parity effects can make this estimate unreliable.

Optional full pair diagnostics use the unnormalized singlet operator
\begin{equation}
 \Delta_{ab}=c_{a\uparrow}c_{b\downarrow}
             -c_{a\downarrow}c_{b\uparrow}
\end{equation}
and retain sign-resolved matrices for rung, leg-0, and leg-1 bonds.  Spatial
decay and sign structure must be reported; a large local value or finite-MPS
pair field alone does not establish superconducting long-range order.

\subsection{Fixed-sector gaps}

Separate QN-conserving DMRG calculations evaluate
\begin{align}
 \Delta_s &= E(N,2)-E(N,0),\\
 \Delta_c &= \tfrac12\left[E(N+2,0)+E(N-2,0)-2E(N,0)\right],\\
 \Ep^{(h)} &= E(N-2,0)+E(N,0)-2E(N-1,1),\\
 \Ep^{(p)} &= E(N+2,0)+E(N,0)-2E(N+1,1),
\end{align}
where the second argument is \(2S_z\).  The present implementation assumes an
even reference \(N\) and an \(S_z=0\) reference sector.  Per-sector sweep
energies, discarded weights, and maximum link dimensions are stored as
convergence evidence.

\section{Numerical acceptance and representative controls}

The following values describe the checked-in Phase-1 base configuration and
are not universal constants.  Every run stores its actual values.
\begin{center}
\begin{tabularx}{\textwidth}{lX}
\toprule
Control & Representative value / rule \\
\midrule
DMRG & 12 sweeps, \(\chi=200\), cutoff \(10^{-10}\), sweep-energy tolerance
  \(10^{-8}\), Krylov dimension 8. \\
Density search & \(|n-n_\star|\le2\times10^{-4}\), at most 16 evaluations,
  initial bracket step \(0.05\), growth factor 2. \\
Fixed-point fields & \(\epsilon_\infty\le10^{-6}\) or
  \(\epsilon_2\le5\times10^{-3}\), including the slow-mode extrapolation, for
  two stable iterations. \\
Energy and identities & corrected energy change per site \(\le10^{-7}\),
  bare-Hamiltonian identity error per site \(\le10^{-9}\), and effective
  eigenvalue error per site \(\le10^{-6}\). \\
Raw orbit probe & 20 raw updates; inspect periods through 2 in the probe;
  periods 1 and 2 enabled; three recurrence links per phase; recurrence
  distance \(\le2\times10^{-6}\) absolute or \(\le10^{-2}\) relative. \\
Period-two shape & worst step cosine \(\le-0.5\) and worst two-step/one-step
  ratio \(\le0.5\). \\
Anderson mixing & memory 5, regularization \(10^{-10}\), initial damping 0.5,
  adaptive range 0.05--0.8. \\
\bottomrule
\end{tabularx}
\end{center}

An accepted solution is therefore a finite-\(L\), finite-\(\chi\) numerical
object.  Phase claims require, at minimum, bond-dimension and length scaling,
boundary/bulk-window checks, consistent order diagnostics, sector gaps, and
seed/basin robustness.  The implementation currently excludes complex fields
and flux order.

\section{Artifacts, provenance, and branch selection}

Each HDF5 state records the complete configuration, applied and measured field
history, correlators and energy decomposition, DMRG sweep evidence, model and
numerical fingerprints, implementation and registry hashes, software/device
metadata, random seed, branch labels, and hash-pinned parent or resume lineage.
Rolling latest/best checkpoints are replaceable; final states and orbit
artifacts are immutable.

Full production states, including MPS objects, live in scratch storage.
Compact stateless mirrors omit MPS tensors but retain the analysis and
selection data plus hashes and the full-artifact location.  A compact state may
be plotted, compared, or used for field-only inheritance, but cannot be used as
an MPS restart.

The comparison tool admits only accepted period-one states or explicitly
enabled, raw-map-validated orbits with matching fingerprints.  For scans, keep
independent, forward, reverse, and metastable continuations as distinct
lineages rather than overwriting a higher-energy basin.

\section{Code map and maintenance}

\begin{center}
\begin{tabularx}{\textwidth}{p{0.34\textwidth}X}
\toprule
File & Responsibility \\
\midrule
\path{src/MeanField.jl} & site ordering, seeds, bare/MF MPO, correlators, and
  geometry-dependent raw field map \\
\path{src/Geometry.jl} & site/rung indexing, supported transverse geometries,
  and spin-resolved density kernels \\
\path{src/Solver.jl} & DMRG schedules, density targeting, SCF driver, and
  checkpoint/diagnostic orchestration \\
\path{src/Convergence.jl} & fixed-point, slow-mode, recurrence, orbit, and
  failure-status gates \\
\path{src/Mixing.jl} & field vectorization, residual metric, linear mixing,
  and Anderson acceleration \\
\path{src/Variational.jl} & channel constants, direct/reconstructed canonical
  energy, and consistency identities \\
\path{src/Diagnostics.jl} & structure factors, \(K_\rho\), entanglement,
  central charge, pair matrices, and fixed-sector gaps \\
\path{src/EpRegistry.jl} & exact/interpolated pair-binding selection and
  weak-coupling checks \\
\path{src/Storage.jl}, \path{src/Provenance.jl}, \path{src/Selection.jl}
  & immutable artifacts, fingerprints, lineage, compact mirrors, and accepted
  branch ranking \\
\bottomrule
\end{tabularx}
\end{center}

When the scientific map changes, update this file together with the relevant
source and tests.  In particular, any change to the \(\alpha/\beta/\nu\) map
must be reconciled with \(E_\perp\), the double-counting correction, and the
Hamiltonian-identity tests.  Campaign outcomes and command histories remain in
\path{docs/RUN_LOG.md}; these notes should contain only durable methods and
clearly labeled representative settings.

\section*{Revision log}
\addcontentsline{toc}{section}{Revision log}

\begin{tabularx}{\textwidth}{llX}
\toprule
Date & Revision & Summary \\
\midrule
2026-09-01 & 0.1 & Initial synthesis of the modular model, SCF/DMRG method,
variational functional, convergence gates, diagnostics, and provenance rules.\\
\bottomrule
\end{tabularx}

\begin{thebibliography}{9}

\bibitem{bollmark2025}
D. Bollmark, N. Kohler, and A. Kantian,
\emph{Phys. Rev. B} \textbf{111}, 125141 (2025),
\href{https://doi.org/10.1103/PhysRevB.111.125141}
{doi:10.1103/PhysRevB.111.125141}.

\bibitem{anderson1965}
D. G. Anderson, ``Iterative procedures for nonlinear integral equations,''
\emph{J. ACM} \textbf{12}, 547--560 (1965).

\bibitem{calabrese2004}
P. Calabrese and J. Cardy, ``Entanglement entropy and quantum field theory,''
\emph{J. Stat. Mech.} P06002 (2004).

\end{thebibliography}

\end{document}
